\section{Tipi di Requisito}
\label{section:req-types}
I requisiti si possono classificare in base a:
\begin{itemize}
    \item \textbf{destinatari}
    \begin{itemize}
        \item \textbf{utente}, si rivolgono direttamente all'utente, usano un linguaggio naturale e diagrammi, astrazione di alto livello
        \item \textbf{di sistema}, si rivolgono principalmente ai progettisti, dettaglio di alto livello
    \end{itemize}
    I requisiti vengono inizialmente concepiti come requisiti utente, così da poter essere valutati da stakeholder meno competenti o tecnici, vengono poi successivamente derivati i corrispettivi requisiti di sistema. Entrambi i tipi di requisito riguardano comportamenti strettamente osservabili e non corrispondono a scelte strettamente di implementazione.
    \item \textbf{tipologia}
    \begin{itemize}
        \item \textbf{funzionali}, indicano che servizi il sistema dovrebbe fornire, come dovrebbe reagire a determinati input e come dovrebbe comportarsi in generale. In generale rientrano in questa categoria tutti i requisiti che indicano funzionalità del sistema
        \item \textbf{qualitativi (o non-funzionali o qualità)}, definiscono proprietà specifiche del sistema misurate secondo precise metriche
        \item \textbf{vincoli}, requisiti che restringono, per ragioni organizzative o tecnologiche, il modo in cui il sistema deve essere sviluppato
        \item \textbf{di dominio}, riguardano da requisiti derivati direttamente dallo specifico dominio o contesto applicativo
    \end{itemize}
\end{itemize}
